Dieses Kapitel beschreibt die vollständige Einrichtung der DropIT-Umgebung.
Die Einrichtung gliedert sich in drei Bereiche: die Registrierung der
Anwendungen im Microsoft Entra Admin Center, die Bereitstellung der
erforderlichen Azure-Dienste sowie den abschließenden Start aller Services
über Docker Compose.
 
% ------------------------------------------------------------
\section{Registrierung im Microsoft Entra Admin Center}
% ------------------------------------------------------------
 
Damit sich Benutzer über Microsoft Single Sign-On anmelden können und die
Backend-Services sicher mit der Microsoft-Infrastruktur kommunizieren,
müssen im Entra Admin Center insgesamt vier App-Registrierungen angelegt
werden: eine für den Client (\texttt{DropIT.Core}) sowie jeweils eine für
die drei Backend-APIs.
 
\subsection{DropIT.Core}
 
Im Entra Admin Center zu \textbf{Identität $\to$ Anwendungen $\to$
App-Registrierungen} navigieren und auf \textbf{+ Neue Registrierung}
klicken. Als Name \texttt{DropIT.Core} eintragen, als unterstützte
Kontotypen \textbf{Nur Konten in diesem Organisationsverzeichnis
(Nur ein Mandant -- DropIT)} wählen und mit \textbf{Registrieren} bestätigen.
 
Anschließend wird ein Client Secret angelegt. Dazu unter
\textbf{Zertifikate \& Geheimnisse $\to$ Clientgeheimnisse} auf
\textbf{+ Neuer geheimer Clientschlüssel} klicken, eine Beschreibung
(z.\,B.\ \texttt{Entwicklung-Key}) sowie ein Ablaufdatum vergeben und mit
\textbf{Hinzufügen} bestätigen. \textbf{Der angezeigte Wert muss sofort
kopiert werden, da er nur einmalig angezeigt wird und später nicht mehr
abrufbar ist.}
 
Zuletzt werden die API-Berechtigungen konfiguriert. Unter
\textbf{API-Berechtigungen $\to$ + Berechtigung hinzufügen $\to$
Von meiner Organisation verwendete APIs} werden nacheinander folgende
Scopes hinzugefügt:
 
\begin{itemize}
    \item \texttt{DropIT.Documents.Api} -- Scope: \texttt{Documents}
    \item \texttt{DropIT.Search.Api} -- Scope: \texttt{Search}
    \item \texttt{DropIT.Zones.Api} -- Scope: \texttt{Zones}
\end{itemize}
 
\subsection{Backend-APIs}
 
Für die drei Backend-Services \texttt{DropIT.Documents.API},
\texttt{DropIT.Zones.API} und \texttt{DropIT.Search.API} wird das
folgende Vorgehen jeweils identisch wiederholt.
 
Unter \textbf{App-Registrierungen} auf \textbf{+ Neue Registrierung}
klicken. Als Name das Schema \texttt{DropIT.<name>.API} verwenden,
als Kontotyp \textbf{Nur ein Mandant (DropIT)} wählen und die
Umleitungs-URI leer lassen. Mit \textbf{Registrieren} bestätigen.
 
Anschließend unter \textbf{Eine API verfügbar machen $\to$
+ Bereich hinzufügen} den Bereichsnamen \texttt{access\_as\_user}
eintragen, die Einwilligung für \textbf{Administratoren und Benutzer}
freigeben und den Status auf \textbf{Aktiviert} setzen.
 
Zusätzlich benötigen zwei der APIs spezifische Graph-Berechtigungen,
die unter \textbf{API-Berechtigungen} hinzugefügt werden:
 
\begin{itemize}
    \item \texttt{DropIT.Zones.API}: \texttt{Files.ReadWrite}, \texttt{Sites.Read.All}
    \item \texttt{DropIT.Documents.API}: \texttt{Sites.Read.All}, \texttt{Sites.ReadWrite}
\end{itemize}
 
% ------------------------------------------------------------
\section{Azure-Ressourcen einrichten}
% ------------------------------------------------------------
 
\subsection{Ressourcengruppe}
 
Im Azure Portal das Menü \textbf{Ressourcengruppen} öffnen und auf
\textbf{+ Erstellen} klicken. Als Ressourcengruppenname empfiehlt sich
ein einheitliches Schema wie \texttt{rg-dropit-core-dev}. Als Region
einen geeigneten Standort wählen (z.\,B.\ \texttt{Germany West Central})
und mit \textbf{Überprüfen + erstellen} und \textbf{Erstellen} abschließen.
 
\subsection{Azure AI Search}
 
Über \textbf{Ressource erstellen} nach \textbf{Azure AI Search} suchen
und den Dienst auswählen. Die zuvor angelegte Ressourcengruppe auswählen
und dieselbe Region wie die übrigen Services verwenden, um Latenzen zu
minimieren. \textbf{Der Dienstname muss global eindeutig sein}, da er
Teil der öffentlichen URL wird. Den Tarif entsprechend des Bedarfs wählen
und mit \textbf{Überprüfen + erstellen} und \textbf{Erstellen} bestätigen.
Falls die gewählte Region keine freien Kapazitäten meldet, kann auf eine
alternative Region ausgewichen werden.
 
% ------------------------------------------------------------
\section{Start mit Docker Compose}
% ------------------------------------------------------------
 
Für den vollständigen Betrieb aller DropIT-Services steht eine fertige
\texttt{docker-compose.yml} bereit, die das Frontend, die Zones-API,
die Documents-API, die Search-API, den AI-Service sowie die
PostgreSQL-Datenbank gemeinsam startet.
 
Vor dem ersten Start müssen alle Zugangsdaten in einer \texttt{.env}-Datei
im Wurzelverzeichnis des Projekts eingetragen werden.
Tabelle~\ref{tab:env-variablen} listet alle erforderlichen Variablen auf.
 
\begin{table}[h]
\centering
\caption{Erforderliche Umgebungsvariablen in der \texttt{.env}-Datei}
\label{tab:env-variablen}
\begin{tabularx}{\textwidth}{lX}
\toprule
\textbf{Variable} & \textbf{Beschreibung} \\
\midrule
\texttt{ENTRA\_TENANT\_ID}          & Tenant-ID des Microsoft Entra Mandanten \\
\texttt{ENTRA\_CLIENT\_ID}          & Client-ID der \texttt{DropIT.Core}-Registrierung \\
\texttt{ENTRA\_CLIENT\_SECRET}      & Das zuvor erstellte Client Secret \\
\texttt{DOCUMENTS\_API\_CLIENT\_ID} & Client-ID der \texttt{DropIT.Documents.API}-Registrierung \\
\texttt{ZONES\_API\_CLIENT\_ID}     & Client-ID der \texttt{DropIT.Zones.API}-Registrierung \\
\texttt{SEARCH\_API\_CLIENT\_ID}    & Client-ID der \texttt{DropIT.Search.API}-Registrierung \\
\texttt{AZURE\_SEARCH\_ENDPOINT}    & URL des Azure AI Search Dienstes \\
\texttt{AZURE\_SEARCH\_KEY}         & Admin-Schlüssel des Azure AI Search Dienstes \\
\texttt{SHAREPOINT\_DRIVE\_ID}      & Drive-ID der SharePoint-Dokumentbibliothek \\
\texttt{OPENAI\_API\_KEY}           & API-Schlüssel für die Dokumentklassifikation \\
\texttt{POSTGRES\_PASSWORD}         & Passwort für die PostgreSQL-Datenbank \\
\texttt{SEARCH\_INTERNAL\_API\_KEY} & Interner Schlüssel für die serviceübergreifende Indexierung \\
\bottomrule
\end{tabularx}
\end{table}
 
Sind alle Werte eingetragen, genügt folgender Befehl, um die gesamte
Anwendung zu starten:
 
\begin{lstlisting}[language=bash, caption={Start aller Services via Docker Compose}]
docker compose up --build
\end{lstlisting}
 
Alle Services starten dann automatisch in isolierten Containern, die
Datenbank wird initialisiert, und das Frontend ist anschließend im
Browser erreichbar.